% 全文排版规则；修改这里后运行 compile.ps1 即可生效。
\usepackage[left=15mm,right=15mm,top=16mm,bottom=17mm,headheight=13pt,headsep=5mm,footskip=8mm]{geometry}
\usepackage{fontspec,setspace}
\usepackage[table]{xcolor}
\IfFontExistsTF{Times New Roman}{\setmainfont{Times New Roman}}{\setmainfont{TeX Gyre Termes}}
\IfFontExistsTF{Arial}{\setsansfont{Arial}}{\setsansfont{TeX Gyre Heros}}
\setmonofont{JetBrainsMono-Light.ttf}[Path=fonts/,Scale=1,BoldFont=JetBrainsMono-Regular.ttf,ItalicFont=JetBrainsMono-LightItalic.ttf,BoldItalicFont=JetBrainsMono-Italic.ttf,RawFeature={-calt}]
\IfFontExistsTF{SimSun}{
  \setCJKmainfont{SimSun}[BoldFont=SimHei,ItalicFont=KaiTi]
  \newCJKfontfamily\tableheaderfont{SimSun}[AutoFakeBold=2]
}{
  \setCJKmainfont{Noto Serif CJK SC}[BoldFont=Noto Sans CJK SC,ItalicFont=FandolKai-Regular.otf]
  \newCJKfontfamily\tableheaderfont{Noto Serif CJK SC}[AutoFakeBold=2]
}
\IfFontExistsTF{STZhongsong}{
  \newCJKfontfamily\headingfont{STZhongsong}[AutoFakeBold=2]
}{
  \newCJKfontfamily\headingfont{Noto Serif CJK SC}[AutoFakeBold=2]
}
\IfFontExistsTF{Segoe UI Symbol}{\newfontfamily\symbolfont{Segoe UI Symbol}}{\newfontfamily\symbolfont{DejaVu Sans}}
% 光学基线校准；保留字号和字宽。设为 0em 可还原。
\newcommand{\CJKBaselineRaise}{0.05em}
\renewcommand{\CJKsymbol}[1]{\raise\CJKBaselineRaise\hbox{#1}}
\renewcommand{\CJKpunctsymbol}[1]{\raise\CJKBaselineRaise\hbox{#1}}
\definecolor{accent}{HTML}{4472C4}
\definecolor{muted}{HTML}{7F7F7F}
\definecolor{tocmuted}{RGB}{128,128,128}
\definecolor{sectionbg}{HTML}{D9E2F3}
\definecolor{headingink}{HTML}{1F3763}
\definecolor{rowshade}{HTML}{F2F2F2}
\definecolor{tablerule}{HTML}{666666}
\newcommand{\ContinuedTableText}{(continued)}
\setstretch{1.03}
\setlength{\parindent}{0pt}
\setlength{\parskip}{3pt plus 1pt minus 1pt}
\hyphenpenalty=500
\exhyphenpenalty=50
\setcounter{secnumdepth}{-2}
\setcounter{tocdepth}{2}
\newcommand{\chapterband}[1]{\begingroup\setlength{\fboxsep}{5pt}\colorbox{accent}{\parbox{\dimexpr\linewidth-2\fboxsep\relax}{\raggedright\color{white}\strut#1\strut}}\endgroup}
\newcommand{\sectionband}[1]{\begingroup\setlength{\fboxsep}{4pt}\colorbox{sectionbg}{\parbox{\dimexpr\linewidth-2\fboxsep\relax}{\raggedright\color{black}\strut#1\strut}}\endgroup}
\newcommand{\subsectionline}[1]{\parbox{\linewidth}{\raggedright\color{headingink}#1\par\nobreak\vspace{1pt}{\color{accent}\hrule height 0.4pt}}}
\ctexset{
  chapter={format=\rmfamily\headingfont\fontsize{16}{21}\selectfont\bfseries,titleformat=\chapterband,beforeskip=0pt,afterskip=3pt,fixskip=true},
  section={format=\headingfont\fontsize{12}{16}\selectfont,titleformat=\sectionband,beforeskip=10pt,afterskip=5pt},
  subsection={format=\rmfamily\headingfont\fontsize{12}{16}\selectfont\bfseries,titleformat=\subsectionline,beforeskip=8pt,afterskip=2pt},
  subsubsection={format=\rmfamily\headingfont\fontsize{11}{14}\selectfont\bfseries,beforeskip=6pt,afterskip=2pt},
  paragraph={format=\rmfamily\headingfont\fontsize{10.5}{13}\selectfont\bfseries,beforeskip=5pt,afterskip=2pt},
  subparagraph={format=\rmfamily\headingfont\fontsize{10.5}{13}\selectfont\bfseries,beforeskip=4pt,afterskip=2pt}
}
\usepackage{longtable,booktabs,array,calc,etoolbox,needspace}
\setlength{\LTpre}{5pt}
\setlength{\LTpost}{5pt}
\setlength{\emergencystretch}{2em}
\AtBeginEnvironment{longtable}{\normalsize\arrayrulecolor{tablerule}\setlength{\heavyrulewidth}{0.65pt}\setlength{\lightrulewidth}{0.5pt}\setlength{\cmidrulewidth}{0.4pt}\setlength{\tabcolsep}{5pt}\setlength{\extrarowheight}{1pt}\setlength{\aboverulesep}{0pt}\setlength{\belowrulesep}{0pt}\renewcommand{\arraystretch}{1.10}}
\usepackage{amsmath,amssymb}
\DeclareMathSizes{10.5}{11}{8}{6}
\usepackage{enumitem}
\setlist{topsep=3pt,itemsep=2pt,parsep=0pt,leftmargin=2.2em}
\providecommand{\tightlist}{\setlength{\itemsep}{1pt}\setlength{\parskip}{0pt}}
\usepackage{fvextra}
\fvset{breaklines=true,breakanywhere=true,fontsize=\normalsize,baselinestretch=1.03}
\AtBeginEnvironment{Verbatim}{\hyphenpenalty=50\exhyphenpenalty=50}
\usepackage{graphicx}
\usepackage{multicol,tocloft}
\setlength{\columnsep}{8mm}
\setlength{\columnseprule}{0.2pt}
\setlength{\cftbeforechapskip}{3pt}
\setlength{\cftbeforesecskip}{0pt}
\setlength{\cftbeforesubsecskip}{0pt}
\renewcommand{\cftchapfont}{\rmfamily\fontsize{12}{14.4}\selectfont\headingfont\bfseries}
\renewcommand{\cftchappagefont}{\rmfamily\fontsize{12}{14.4}\selectfont\bfseries}
\renewcommand{\cftsecfont}{\fontsize{10.5}{12.6}\selectfont}
\renewcommand{\cftsecpagefont}{\fontsize{10.5}{12.6}\selectfont}
\renewcommand{\cftsubsecfont}{\fontsize{9}{10.8}\selectfont\color{tocmuted}}
\renewcommand{\cftsubsecpagefont}{\fontsize{9}{10.8}\selectfont\color{tocmuted}}
\cftsetindents{chapter}{0pt}{0pt}
\cftsetindents{section}{5pt}{0pt}
\cftsetindents{subsection}{10pt}{0pt}
\makeatletter
\newcommand{\manualtoc}{\begingroup\setstretch{1}\hypersetup{linkcolor=black}\setlength{\parskip}{0pt}\@starttoc{toc}\endgroup}
\makeatother
\usepackage{fancyhdr}
\pagestyle{fancy}
\fancyhf{}
\fancyhead[C]{\footnotesize\color{muted}\nouppercase{\leftmark}}
\fancyfoot[LE,RO]{\normalsize\thepage}
\renewcommand{\headrulewidth}{0pt}
\renewcommand{\footrulewidth}{0pt}
\renewcommand{\chaptermark}[1]{\markboth{#1}{}}
\fancypagestyle{plain}{\fancyhf{}\fancyfoot[LE,RO]{\normalsize\thepage}\renewcommand{\headrulewidth}{0pt}}
\usepackage{xurl}
\usepackage[unicode,colorlinks=true,linkcolor=accent,urlcolor=accent,bookmarksopen=true,bookmarksnumbered=false]{hyperref}
\hypersetup{pdflang={en-US},pdftitle={Annotation Reviewer User Manual},pdfsubject={Operating instructions and data conventions}}
\usepackage{bookmark}
% PDF 书签：章为第一级，节为第二级，小节为第三级。
\bookmarksetup{addtohook={%
  \bookmarksetup{bold=false,color=black}%
  \ifnum\bookmarkget{level}=0 \bookmarksetup{bold}\fi
  \ifnum\bookmarkget{level}=2 \bookmarksetup{color=tocmuted}\fi
}}
% Contents-only layout, independent of regenerated Markdown content.
\setlength{\cftbeforechapskip}{8pt}
\bookmarksetup{open=true,openlevel=1}
\makeatletter
\newcount\manual@toclevel
\def\manual@splitnumber#1 #2\manual@stop{%
  \def\manual@tocnumber{#1}%
  \def\manual@toctitle{#2}%
}
\newcommand{\manual@numberedcontentsline}[4]{%
  \def\manual@kind{#1}%
  \def\manual@chapter{chapter}%
  \def\manual@section{section}%
  \ifx\manual@kind\manual@chapter
    \manual@toclevel=0
  \else\ifx\manual@kind\manual@section
    \manual@toclevel=1
  \else
    \manual@toclevel=2
  \fi\fi
  \begingroup
    \manual@splitnumber#2\manual@stop
    \ifcase\manual@toclevel
      \settowidth{\cftchapnumwidth}{\cftchapfont\manual@tocnumber\enspace}%
    \or
      \settowidth{\cftsecnumwidth}{\cftsecfont\manual@tocnumber\enspace}%
    \else
      \settowidth{\cftsubsecnumwidth}{\cftsubsecfont\manual@tocnumber\enspace}%
    \fi
    \ifnum\manual@toclevel=2
      \color{tocmuted}\hypersetup{linkcolor=tocmuted}%
    \else
      \color{black}\hypersetup{linkcolor=black}%
    \fi
    \manual@savedcontentsline{#1}{\numberline{\manual@entrystrut\manual@tocnumber}\manual@toctitle}{#3}{#4}%
  \endgroup
}
% Choose column breaks locally; a third-level entry may start a column.
\newcount\manual@total
\newcount\manual@start
\newcount\manual@index
\newcount\manual@natural
\newcount\manual@break
\newcount\manual@candidate
\newcount\manual@prior
\newcount\manual@rank
\newcount\manual@rows
\newcount\manual@shift
\newcount\manual@column
\newdimen\manual@capacity
\newdimen\manual@used
\newdimen\manual@fullheight
\newdimen\manual@remaining
\newdimen\manual@needed
\newdimen\manual@before
\newdimen\manual@strutheight
\newdimen\manual@strutdepth
\newif\ifmanual@legal
\newcommand{\manual@entrystrut}{%
  \vrule width0pt height\manual@strutheight depth\manual@strutdepth}
\def\manual@level#1{\csname manual@level@\number#1\endcsname}
\def\manual@lines#1{\csname manual@lines@\number#1\endcsname}
\def\manual@height#1{\dimexpr
  \ht\csname manual@box@\number#1\endcsname+
  \dp\csname manual@box@\number#1\endcsname\relax}
\def\manual@toc@chapter{0}
\def\manual@toc@section{1}
\def\manual@toc@subsection{2}
\newcommand{\manual@contentsline}[4]{%
  \ifcsname manual@toc@#1\endcsname
    \advance\manual@total\@ne
    \expandafter\newbox\csname manual@box@\the\manual@total\endcsname
    \expandafter\edef\csname manual@level@\the\manual@total\endcsname{%
      \csname manual@toc@#1\endcsname}%
    \setbox\csname manual@box@\the\manual@total\endcsname=\vbox{%
      \hsize=\linewidth \parskip=0pt
      \manual@numberedcontentsline{#1}{#2}{#3}{#4}%
      \expandafter\xdef\csname manual@lines@\the\manual@total\endcsname{%
        \ifnum\prevgraf>0 \the\prevgraf\else 1\fi}%
    }%
    % Give the last line the same minimum depth as an ordinary single-spaced row.
    \ifdim\dp\csname manual@box@\the\manual@total\endcsname<\manual@strutdepth
      \dp\csname manual@box@\the\manual@total\endcsname=\manual@strutdepth
    \fi
  \fi
}
\newcommand{\manual@checkbreak}[1]{%
  \manual@legaltrue
  \manual@prior=#1\relax \advance\manual@prior\m@ne
  \ifnum\manual@prior<\manual@start\else
    % A chapter stays with its next entry; a section with its first subsection.
    \ifnum\manual@level\manual@prior=0
      \manual@legalfalse
    \else\ifnum\manual@level\manual@prior=1
      \ifnum\manual@level#1=2 \manual@legalfalse\fi
    \fi\fi
  \fi
}
\newcommand{\manual@validbreak}{%
  \manual@checkbreak\manual@break
  \ifmanual@legal\else
    \advance\manual@break\m@ne
    \expandafter\manual@validbreak
  \fi
}
\newcommand{\manual@preferheading}{%
  \manual@candidate=\manual@break
  \manual@before=\manual@used
  \manual@rank=3 \manual@rows=0 \manual@shift=0
  \@whilenum\manual@candidate>\manual@start\do{%
    \manual@checkbreak\manual@candidate
    \ifdim\manual@before<\manual@needed \manual@legalfalse\fi
    \ifmanual@legal
      \ifnum\manual@level\manual@candidate=0
        \ifnum\manual@rank>0
          \manual@break=\manual@candidate \manual@rank=0 \manual@shift=\manual@rows
        \fi
      \else\ifnum\manual@level\manual@candidate=1
        \ifnum\manual@rows<3
          \ifnum\manual@rank>1
            \manual@break=\manual@candidate \manual@rank=1 \manual@shift=\manual@rows
          \fi
        \fi
      \fi\fi
    \fi
    \advance\manual@candidate\m@ne
    \advance\manual@rows\manual@lines\manual@candidate\relax
    \advance\manual@before-\manual@height\manual@candidate
    \ifnum\manual@rows>3 \manual@candidate=\manual@start\fi
  }%
}
\newcommand{\manual@fillcolumn}{%
  \manual@used=0pt \manual@index=\manual@start
  \manual@natural=\numexpr\manual@start-1\relax
  \@whilenum\manual@index<\numexpr\manual@total+1\relax\do{%
    \ifdim\dimexpr\manual@used+\manual@height\manual@index\relax>\manual@capacity
      \manual@index=\numexpr\manual@total+1\relax
    \else
      \advance\manual@used\manual@height\manual@index
      \manual@natural=\manual@index \advance\manual@index\@ne
    \fi
  }%
  \manual@break=\numexpr\manual@natural+1\relax
  \ifnum\manual@natural<\manual@total
    \manual@validbreak
    \manual@natural=\numexpr\manual@break-1\relax
    \manual@used=0pt \manual@index=\manual@start
    \@whilenum\manual@index<\manual@break\do{%
      \advance\manual@used\manual@height\manual@index
      \advance\manual@index\@ne
    }%
  \fi
}
\newcommand{\manual@layouttoc}{%
  \manual@start=1 \manual@column=1
  \@whilenum\manual@start<\numexpr\manual@total+1\relax\do{%
    \manual@fullheight=\@colroom
    \ifnum\manual@column>2 \manual@fullheight=\textheight\fi
    \manual@capacity=\manual@fullheight \manual@needed=-\maxdimen
    % Balance the final pair of columns without creating an extra page.
    \ifodd\manual@column
      \manual@remaining=0pt \manual@index=\manual@start
      \@whilenum\manual@index<\numexpr\manual@total+1\relax\do{%
        \advance\manual@remaining\manual@height\manual@index
        \advance\manual@index\@ne
      }%
      \ifdim\manual@remaining>2\manual@fullheight\else
        \manual@needed=\dimexpr\manual@remaining-\manual@fullheight\relax
        \manual@capacity=.5\manual@remaining
        \advance\manual@capacity .5\baselineskip
      \fi
    \fi
    \manual@fillcolumn
    \ifdim\manual@used<\manual@needed
      \manual@capacity=\manual@fullheight \manual@needed=-\maxdimen
      \manual@fillcolumn
    \fi
    \manual@shift=0
    \ifnum\manual@natural<\manual@total
      \manual@preferheading
    \fi
    \typeout{MANUALTOC column=\the\manual@column;
      start=\the\manual@start; natural=\the\manual@natural;
      next=\the\manual@break; shift=\the\manual@shift;
      capacity=\the\manual@capacity}%
    \manual@index=\manual@start
    \@whilenum\manual@index<\manual@break\do{%
      \nointerlineskip
      \box\csname manual@box@\the\manual@index\endcsname
      \penalty\@M
      \advance\manual@index\@ne
    }%
    \manual@start=\manual@break
    \ifnum\manual@start<\numexpr\manual@total+1\relax
      \columnbreak
      \advance\manual@column\@ne
    \fi
  }%
}
\renewcommand{\manualtoc}{%
  \begingroup
    \setstretch{1}\setlength{\parskip}{0pt}%
    \manual@strutheight=\ht\strutbox \manual@strutdepth=\dp\strutbox
    \manual@total=0
    \let\manual@savedcontentsline\contentsline
    \let\contentsline\manual@contentsline
    \raggedcolumns
    \begingroup
      \makeatletter
      \@input{\jobname.toc}%
      \manual@layouttoc
      \if@filesw
        \expandafter\newwrite\csname tf@toc\endcsname
        \immediate\openout\csname tf@toc\endcsname\jobname.toc\relax
      \fi
      \@nobreakfalse
    \endgroup
  \endgroup
}
\makeatother

\widowpenalty=10000
\clubpenalty=10000
\raggedbottom
\input{build-info.tex}
\IfFileExists{tweaks.tex}{% 持久微调：重新生成不会覆盖本文件。
% 示例：统一调整颜色、正文行距、表格上下间距。
% \definecolor{accent}{HTML}{315F86}
% \setstretch{1.03}
% \setlength{\LTpre}{5pt}
% 中文视觉高度微调：修改下面这一行后运行“仅编译”；0.00em 表示不抬高，0.01em 表示中文抬高 0.01em。
\renewcommand{\CJKBaselineRaise}{-0.02em}
}{}
